\addtocontents{toc}{\protect\setcounter{tocdepth}{-1}}
\setcounter{chapter}{1}
\chapter{Background and Related Work}

\chapter{General Trust Assessment Methodology }
This chapter introduces the \emph{General Trust Assessment Methodology} that forms the conceptual backbone of this dissertation. 
The methodology builds upon earlier work on trust evaluation developed in the context of connected and autonomous systems, including the Evidence-based \ac{TAF} developed within the Horizon Europe CONNECT project~\cite{connect_d31_taf}, the 5GAA white paper on trust evaluation~\cite{5gaa_trust4cav_2025}, the trust assessment framework presented at the International Conference on Information and Communications Security (ICICS)~\cite{fotos2025actions}, and subsequent work on trustworthy data assessment presented at the Bias Workshop at ECML~\cite{ecmldatatrust}. 
In this chapter, these contributions are unified and generalized into a systematic methodology for structuring trust assessment problems and deriving trust evaluations so that it can be instantiated across different \ac{AI} contexts.



\chapter{Optimization of \acs{DSPG} Synthesis and \aclu{TNA-SL}}
This chapter presents an optimized framework for \acf{TNA-SL} which is often used to generate expression in phase 1 of the Trust Assessment methodology. The chapter introduces foundational, algorithmic, and operator level improvements that address several limitations of the classical \ac{TNA-SL} methodology. In particular, we propose a refined structural representation of trust networks based on \acf{PNPS}, derive optimized \ac{DSPG} synthesis criteria that preserve more information, develop an improved \ac{DSPG} analysis algorithm that reduces computational overhead, and introduce a new referral-edge path discounting operator ensuring consistent trust propagation over referral chains. The \acf{PNPS} structural concept was initially introduced through collaboration with colleagues at Huawei and research partners in Luxembourg. Building on this definition, all theoretical developments presented in this chapter, including the formalization, analysis, and proofs of the proposed properties and algorithms, were carried out in this work.

The results presented in this chapter are based on two peer-reviewed publications. The structural and algorithmic optimization of \ac{DSPG} synthesis and analysis were published in \cite{10.1007/978-3-032-08887-1_4} and presented at the RuleML+RR 2025 conference. The referral-edge trust discounting operator and its theoretical properties were presented in \cite{10706345} at the FUSION 2024 conference. This chapter unifies these contributions into a coherent framework for efficient and information-preserving trust inference.

\chapter{\acf{PaTAS}}
The previous chapters established the conceptual and methodological foundations for trust assessment in general and particularly in \ac{AI} systems. In particular, \cref{chap:Trust_method} introduced a general evidence-based methodology for trust evaluation, while \cref{chap:tnasl} improved the graph-based reasoning mechanisms required to compute trust expressions. Building on these foundations, this chapter introduces the \emph{\acf{PaTAS}}, a system architecture designed to operationalize the proposed methodology in \ac{ML} pipelines. \ac{PaTAS} provides a scalable and modular framework for performing trust assessment throughout both training and inference, starting from the training dataset. The system enables the evaluation of \ac{NN} outputs in parallel with their inference by allowing different trust sources and reasoning components to be processed independently, without impacting the standard \ac{NN} inference process. This design allows efficient evaluation of trust in complex \ac{NN} systems where multiple evidence streams must be combined. The concepts presented in this chapter build upon previous publication. In particular, the architecture and core reasoning principles are derived from the work presented in~\cite{patas}, currently under review at the \emph{Elsevier Knowledge Based Systems}.

\chapter{Definition and Assessment of AI Training Datasets Trustworthiness}
\label{chap:data}
The previous chapter presented the overall architecture of \ac{PaTAS}, highlighting its two main components: (i) the Trust Assessment Module focusing on the assessment of the
trustworthiness of the training dataset and the input of the \ac{NN}, and (ii) the Trust Propagation Module focusing on the propagation of trust through the model, from the training data to the learned parameters and from the input to the output via these parameters. This chapter focuses on the first component, namely the trust assessment module. The underlying ideas were initially explored in paper presented at the Bias Workshop at ECML~\cite{ecmldatatrust} and further discussed in the position paper~\cite{trust-aidatatrust}. In this dissertation, these contributions are unified and extended to provide a comprehensive and formal framework for assessing the trustworthiness of training datasets and input data of a \ac{NN}.

\chapter{Trust Propagation in Neural Network}
The previous \cref{chap:ovl} presented the overall architecture of \ac{PaTAS}, highlighting its two main components: (i) the assessment of trust in the data, and (ii) the propagation of trust through the model, from the training data to the learned parameters and from the input to the output via these parameters. While the trustworthiness of the data was addressed in the preceding chapter, this chapter focuses on the second component, namely the propagation of trust within \aclu{NN}. In particular, it details how trust, initially derived from the dataset, is integrated into the model parameters during training and subsequently propagated from the input to the output during inference. This chapter formalizes the mechanisms that enable consistent and interpretable trust propagation across the different stages of the learning process, providing the operational core of the \ac{PaTAS} framework.

\chapter{Use Case: Applying the Trust Assessment Methodology to ...}
\noindent
The previous chapters introduced the theoretical foundations and system architecture required to perform trust assessment in \ac{NN} systems. In particular, \cref{chap:Trust_method} presented a general methodology for evidence-based trust evaluation, \cref{chap:tnasl} developed the graph-based reasoning framework for trust expression evaluation, and \cref{chap:ovl,chap:data,chap:patas} introduced the \ac{PaTAS} architecture for scalable trust assessment through the whole \ac{NN} system pipeline. This chapter demonstrates how these can be applied in practice through a concrete use case. The objective is not to introduce new theoretical contributions, but rather to illustrate how the proposed methodology can be applied to assess trustworthiness properties within a realistic \ac{ML} scenario. By instantiating the full trust assessment pipeline, from trust proposition definition to evidence collection, opinion quantification, and reasoning, we show how the framework can support practical decision-making regarding the
trustworthiness of \ac{ML} systems. The use case highlights how different trust sources can be integrated, how evidence is translated into subjective opinions, and how the resulting trust expressions can be evaluated within the proposed framework. This demonstration provides an end-to-end view of the trust assessment process and illustrates the practical relevance of the work developed throughout this dissertation.

\chapter{Calibration-Based Trust Assessment of Neural Networks}
\noindent
The previous chapters introduced and evaluated a framework for assessing the trustworthiness of \ac{NN} systems through its pipeline based on its internal behaviour. However, in many practical settings \ac{ML} models are deployed as black boxes, making it difficult to assess trust based on internal components or training processes. In such cases, trust must be inferred from observable properties of the model. One important property is calibration, which measures the alignment between predicted probabilities and empirical outcomes. In this chapter, we propose a method to quantify \ac{NN} calibration within the Evidence based \aclu{TAF} introduced in this dissertation. This chapter builds on work published and presented at the International Conference on Information Fusion (Fusion 2025)\cite{11124121}.

